\section{Policy Learning：Policy Gradient 方法}

为了克服 Q-Learning 的局限，本节介绍另一类 Deep RL 算法------\textbf{Policy Gradient}（策略梯度），它直接学习策略函数 $\pi(s)$，无需学习中间的 Q-function。

\subsection{Policy Gradient 的核心思想}

\begin{figure}[H]
\centering
\includegraphics[width=0.85\textwidth]{slides-images/slide-38.jpg}
\caption{离散动作空间中的策略：输出每个动作的概率分布 $P(a|s)$\protect\footnotemark}
\end{figure}
\footnotetext{Slides 第38页。}

与 DQN 不同，Policy Gradient 方法不再估计每个动作的 Q 值，而是直接输出一个\textbf{动作上的概率分布}。网络的输入是状态 $s$，输出是 $P(a|s)$------在该状态下采取每个动作的概率。

对于\textbf{离散}动作空间，输出是一个 softmax 概率分布（类似分类任务）：

$$P(a = \text{left} | s) = 0.6, \quad P(a = \text{stay} | s) = 0.25, \quad P(a = \text{right} | s) = 0.15$$

\subsection{处理连续动作空间}

Policy Gradient 最大的优势在于能够自然地处理\textbf{连续动作空间}。

\begin{figure}[H]
\centering
\includegraphics[width=0.95\textwidth]{slides-images/slide-40.jpg}
\caption{Policy Gradient 处理连续动作空间：网络输出高斯分布的参数 $(\mu, \sigma^2)$\protect\footnotemark}
\end{figure}
\footnotetext{Slides 第40页。}

对于连续动作空间，网络不再输出各动作的概率，而是输出一个\textbf{概率分布的参数}。例如，假设动作服从高斯分布：

$$P(a|s) = \mathcal{N}(\mu, \sigma^2)$$

网络输出均值 $\mu$ 和方差 $\sigma^2$，然后从该分布中采样得到具体动作：

$$\pi(s) \sim P(a|s) = \mathcal{N}(\mu, \sigma^2)$$

\begin{importantbox}{Policy Gradient 的关键优势}
\begin{itemize}
    \item \textbf{处理连续动作空间}：通过输出概率分布的参数（如高斯分布的 $\mu$ 和 $\sigma^2$），可以自然地处理连续控制问题
    \item \textbf{学习随机策略}：输出的是概率分布而非确定性动作，天然支持 Stochastic Policy
    \item \textbf{直接优化目标}：直接优化策略以最大化期望回报，无需学习中间的 Q-function
\end{itemize}
\end{importantbox}

\subsection{Policy Gradient 的训练算法}

\begin{figure}[H]
\centering
\includegraphics[width=0.85\textwidth]{slides-images/slide-44.jpg}
\caption{Policy Gradient 的训练流程：五步训练算法\protect\footnotemark}
\end{figure}
\footnotetext{Slides 第44页。}

Policy Gradient 的训练遵循以下五步流程：

\begin{enumerate}
    \item \textbf{初始化 Agent}：随机初始化策略网络的参数
    \item \textbf{运行策略直到终止}：让 Agent 按照当前策略在环境中运行一个完整的 episode
    \item \textbf{记录所有 (state, action, reward)}：收集整个 episode 中的所有轨迹数据
    \item \textbf{降低低奖励动作的概率}：对于导致低回报的动作，减少其被选中的概率
    \item \textbf{提高高奖励动作的概率}：对于导致高回报的动作，增加其被选中的概率
\end{enumerate}

\subsection{Policy Gradient 的损失函数}

训练的数学核心在于定义合适的损失函数：

$$\text{loss} = -\log P(a_t | s_t) \cdot R_t$$

\begin{figure}[H]
\centering
\includegraphics[width=0.95\textwidth]{slides-images/slide-46.jpg}
\caption{Policy Gradient 的损失函数与梯度更新公式\protect\footnotemark}
\end{figure}
\footnotetext{Slides 第46页。Williams Reinforcement Learning 1992.}

\begin{itemize}
    \item $-\log P(a_t | s_t)$：选择动作 $a_t$ 的负对数似然（log-likelihood）
    \item $R_t$：从时刻 $t$ 开始的折扣累积回报
\end{itemize}

梯度下降更新规则为：

$$w' = w - \nabla \text{loss} = w + \nabla \log P(a_t | s_t) \cdot R_t$$

\begin{importantbox}{Policy Gradient 损失函数的直觉理解}
这个损失函数的设计非常精巧：
\begin{itemize}
    \item 当 $R_t > 0$（高回报）时，loss 的负号使得梯度更新\textbf{增大} $P(a_t|s_t)$，即提高该动作的概率
    \item 当 $R_t < 0$（低回报/惩罚）时，梯度更新\textbf{减小} $P(a_t|s_t)$，即降低该动作的概率
    \item $R_t$ 的绝对值大小决定了更新幅度------回报越大，对策略的调整越强
\end{itemize}
这就是 REINFORCE 算法（Williams, 1992）的核心思想。
\end{importantbox}

\begin{warningbox}{Policy Gradient 的高方差问题}
Policy Gradient 方法的一个已知问题是\textbf{高方差}（High Variance）。由于使用 Monte Carlo 采样来估计回报，不同 episode 之间的回报波动可能很大，导致梯度估计噪声较大、训练不稳定。常见的缓解方法包括：使用 baseline 减小方差（如 Actor-Critic 方法）、增加采样数量、使用 reward normalization 等。
\end{warningbox}

\subsection{本章小结}

Policy Gradient 方法直接学习策略函数，通过"提高高回报动作的概率、降低低回报动作的概率"来优化策略。其最大优势是能够处理连续动作空间并学习随机策略。损失函数由动作的对数似然和折扣回报的乘积构成，通过梯度上升来优化期望回报。


